\section*{Background}

A natural disaster can increase firms' need to replace productive assets (e.g. machinery, facilities) while disrupting the revenues available to finance investment. Cash accumulated before the event may help bridge this funding gap. Hurricane Helene's impact on the southeastern United States in September 2024 provides a setting to examine this financing problem.

Duchin, Ozbas and Sensoy (2010) find that corporate investment declined more strongly among firms with lower cash reserves during the financial crisis beginning in 2007. Their findings suggest that firms can use internal funds to maintain investment when external finance becomes harder to obtain. This proposal examines whether a similar relationship arises after a physical shock, which also damages productive assets and changes investment needs.

The research question is: among manufacturing firms domiciled and listed in the United States, is the change in capital expenditure associated with exposure to Helene more positive for firms with higher pre-storm cash holdings? The proposed analysis would compare exposed and unexposed firms before the storm and during the first four full quarters afterwards (2024Q4 to 2025Q3). This comparison would distinguish differences in the response to Helene from a general tendency for firms with more cash to invest more.
